\documentclass[11pt,a4paper]{article}
\usepackage[utf8]{inputenc}
\usepackage{amsmath,amssymb}
\usepackage{geometry}
\usepackage{booktabs}
\usepackage{xcolor}
\usepackage{listings}
\usepackage{hyperref}

\geometry{margin=1in}

\title{\textbf{Day 3: NLP, Encoder-Decoder and Transformers}}
\author{Data Science Foundation Series}
\date{\today}

\begin{document}

\maketitle

\section{Natural Language Processing \& Tokenization}
\begin{itemize}
    \item \textbf{Natural Language Processing (NLP):} The discipline of enabling computational models to parse, represent, and generate natural human language.
    \item \textbf{Tokenization Levels:} Character-level (fine-grained, zero OOV, verbose), Word-level (high semantics, high OOV), and Subword-level (Byte-Pair Encoding, WordPiece) balancing vocabulary size and sequence length.
    \item \textbf{Dense Word Embeddings:} Mapping discrete token indices to continuous feature vectors $\mathbf{e} \in \mathbb{R}^d$ such that semantic proximity corresponds to geometric alignment:
    \begin{equation}
    \text{Cosine Similarity}(\mathbf{u}, \mathbf{v}) = \frac{\mathbf{u} \cdot \mathbf{v}}{\|\mathbf{u}\|_2 \|\mathbf{v}\|_2}
    \end{equation}
\end{itemize}

\section{Attention and Self-Attention Mechanisms}
\subsection{The Query-Key-Value Formulation}
Given input matrix $\mathbf{X} \in \mathbb{R}^{N \times d_{\text{model}}}$, linear projections produce Queries $\mathbf{Q} = \mathbf{X}\mathbf{W}_Q$, Keys $\mathbf{K} = \mathbf{X}\mathbf{W}_K$, and Values $\mathbf{V} = \mathbf{X}\mathbf{W}_V$.

\subsection{Scaled Dot-Product Attention}
\begin{equation}
\text{Attention}(\mathbf{Q}, \mathbf{K}, \mathbf{V}) = \text{softmax}\left(\frac{\mathbf{Q}\mathbf{K}^T}{\sqrt{d_k}}\right)\mathbf{V}
\end{equation}
The scaling factor $\frac{1}{\sqrt{d_k}}$ counteracts the growth of dot product magnitudes with increasing dimensionality $d_k$, preventing the softmax activation from saturating and producing vanishing gradients.

\subsection{Multi-Head Attention (MHA)}
\begin{equation}
\text{MultiHead}(\mathbf{Q}, \mathbf{K}, \mathbf{V}) = \text{Concat}(\text{head}_1, \dots, \text{head}_h)\mathbf{W}^O
\end{equation}
where $\text{head}_i = \text{Attention}(\mathbf{Q}\mathbf{W}_i^Q, \mathbf{K}\mathbf{W}_i^K, \mathbf{V}\mathbf{W}_i^V)$.

\section{Positional Encoding \& Transformer Architecture}
\subsection{Sinusoidal Positional Encoding}
Because self-attention is permutation-invariant, positional coordinates are injected into embeddings:
\begin{align}
PE_{(pos, 2i)} &= \sin\left(\frac{pos}{10000^{2i/d_{\text{model}}}}\right) \\
PE_{(pos, 2i+1)} &= \cos\left(\frac{pos}{10000^{2i/d_{\text{model}}}}\right)
\end{align}

\subsection{Residual Connections and Normalization}
Each sublayer incorporates residual skip connections and layer normalization:
\begin{equation}
\mathbf{y} = \text{LayerNorm}(\mathbf{x} + \text{Sublayer}(\mathbf{x}))
\end{equation}

\section{The Architectural Trinity: BERT, GPT, and T5}
\begin{itemize}
    \item \textbf{BERT (Encoder-Only):} Bidirectional attention mask; trained via Masked Language Modeling (MLM); optimized for sequence classification, NER, and extractive search.
    \item \textbf{GPT (Decoder-Only):} Causal lower-triangular attention mask; trained via autoregressive next-token prediction; optimized for generative text completion and reasoning.
    \item \textbf{T5 (Encoder-Decoder):} Full bidirectional encoder coupled with causal autoregressive decoder via cross-attention; unified text-to-text transformation framework.
\end{itemize}

\section{Autoregressive Generation \& Decoding Strategies}
Given prefix tokens $x_{1:t}$, the model emits unnormalized logits $\mathbf{z}_{t+1}$. Softmax yields next-token probabilities:
\begin{equation}
P(w_i | x_{1:t}) = \frac{\exp(z_i / T)}{\sum_j \exp(z_j / T)}
\end{equation}
where $T$ represents temperature. Decoding strategies include Greedy Search ($\text{argmax}$), Top-$K$ truncation, and Top-$P$ (Nucleus) cumulative probability sampling.

\end{document}
